\documentclass[11pt, a4paper]{article}

\usepackage[utf8]{inputenc}
\usepackage[italian]{babel}
\usepackage{geometry}
\geometry{top=2.5cm, bottom=2.5cm, left=2.5cm, right=2.5cm}

\usepackage{float}
\usepackage{amsmath}
\usepackage{graphicx}
\usepackage{booktabs} % Tabelle di qualità professionale
\usepackage{listings} % Per formattare codice o messaggi del protocollo
\usepackage{xcolor}
\usepackage{hyperref}
\usepackage{bytefield}
\usepackage{parskip}

\usepackage{lipsum}

% Configurazione estetica per i log o codice
\lstset{
	basicstyle=\ttfamily\small,
	backgroundcolor=\color{gray!10},
	frame=single,
	breaklines=true
}

\title{\textbf{\Large Relazione di Progetto di Sistemi Operativi}\\ \huge Applicazione Elenco Telefonico con Socket TCP}
\author{\textbf{Giacomo Abbruciati} - Matricola: 0356240}

\begin{document}
	
	\maketitle
	
	\tableofcontents
	\newpage
	
	\section{Introduzione, Specifiche e Scelte progettuali}
	\lipsum[1][1]

	\subsection{Specifiche del progetto}
	\lipsum[1][1-4]
	
	\subsection{Scelte progettuali}
	\lipsum[2]
	
	\section{Specifica del Protocollo}
	\lipsum[1]
	
	\subsection{Struttura Header di Pacchetto}
	
	L'header di pacchetto contiene le informazioni relative alla natura del messaggio e alla lunghezza del relativo payload.\\	
	
	Il primo byte rappresenta l'\textit{Opcode} (\texttt{0x00}--\texttt{0x03}) della richiesta se il messaggio è in ingresso al server (\text{Client} $\rightarrow$ \text{Server}), oppure lo \textit{Status Code} (\texttt{0x80}--\texttt{0x99}) se il messaggio è in uscita dal server (\text{Server} $\rightarrow$ \text{Client}).
	
	I successivi due byte indicano la lunghezza del \textit{Payload}, ovvero il contenuto effettivo del messaggio.
	
	%TODO cambia range di status code
	
	\begin{figure}[htbp]
		\centering
		\begin{bytefield}[bitwidth=1.3em]{24}
			\bitheader{0,7,8,23} \\
			\bitbox{8}{Opcode / Status code} & \bitbox{16}{Payload Length}
		\end{bytefield}
		\caption{Struttura dell'Header di Pacchetto}
	\end{figure}
	
	\subsection{Opcode di richiesta}

	Il client può inviare comandi al server blablabla
	
	\begin{table}[htbp]
		\centering
		\begin{tabular}{llll}
			\toprule
			\textbf{Hex} & \textbf{Opcode} & \textbf{Parametri} & \textbf{Descrizione} \\
			\midrule
			\texttt{0x00} & \texttt{LOGIN} & \texttt{<username> <password\_hash>} & Effettua l'accesso\\
			\texttt{0x01} &	\texttt{REGISTER} & \texttt{<username> <password>} & Registra un utente\\
			\texttt{0x02} &	\texttt{INSERT} & \texttt{<name> <phone\_number>} & Inserisce un contatto\\
			\texttt{0x03} &	\texttt{SEARCH} & \texttt{<name>} & Ricerca un contatto\\
			\bottomrule
		\end{tabular}
		\caption{Lista di opcode disponibili}
	\end{table}
	
	\subsubsection{Opcode \texttt{LOGIN} \& \texttt{REGISTER}}
	
	Gli \textit{Opcode} \texttt{LOGIN} e \texttt{REGISTER} individuano rispettivamente le operazioni di Accesso e Registrazione al servizio di elenco telefonico. Il \textit{Payload} è strutturato ponendo nei primi $\text{Payload Length} - 32$ byte l'username dell'utente e negli ultimi 32 byte l'hash della password effettuato dal client	.
	
	\begin{figure}[htbp]
		\centering
		\begin{bytefield}[bitwidth=1.1em]{24}
			\wordbox{2}{Username (variabile -- $N$ byte)} \\
			\wordbox{3}{Password Client Hash (32 Byte -- SHA-256)}
		\end{bytefield}
		\caption{Struttura del payload di una richiesta \texttt{LOGIN} o \texttt{REGISTER}.}
		\label{fig:login_payload}
	\end{figure}
	
	\subsection{Status code di risposta}
	
	\begin{table}[H] %cambia in htbp 
	\centering
		\begin{tabular}{lll}
			\toprule
			\textbf{Hex} & \textbf{Status} & \textbf{Descrizione} \\
			\midrule
			\texttt{0x80} & \texttt{OK} & Operazione completata con successo \\
			\texttt{0x81} & \texttt{SERVER\_ERROR} & Errore server \\
			\texttt{0x82} & \texttt{200 OK [dati]} & Operazione completata con successo \\
			\texttt{200 OK [dati]} & \texttt{200 OK [dati]} & Operazione completata con successo \\
			\bottomrule
		\end{tabular}
	\caption{Codici di stato restituiti dal Server}
	\end{table}
	
	\section{Specifica del Server}
	\subsection{Architettura Concorrente}
	Il server è stato progettato per gestire più client contemporaneamente. Per ogni nuova connessione accettata sulla socket di ascolto (porta \texttt{8080}), viene creato un nuovo \textit{thread} dedicato alla gestione della sessione.
	
	\subsection{Struttura Dati e Sincronizzazione}
	I dati dell'elenco telefonico sono mantenuti in memoria tramite una Mappa (Chiave: \texttt{Nome}, Valore: \texttt{Numero}).
	Poiché la risorsa è condivisa tra più thread, gli accessi in scrittura (\texttt{ADD}, \texttt{DELETE}) sono protetti tramite un meccanismo di **Sezione Critica (Mutex)** per evitare \textit{race condition}.
	
	\section{Specifica del Client}
	Il client è un'applicazione da riga di comando che:
	\begin{enumerate}
		\item Apre una socket TCP verso l'IP e la porta del server.
		\item Mostra un menu o legge i comandi dell'utente da \texttt{stdin}.
		\item Formatta la richiesta secondo il protocollo e la invia sulla socket.
		\item Attende la risposta del server, ne fa il parsing e mostra il risultato all'utente.
	\end{enumerate}
	
	\section{Esecuzione e Collaudo}
	\subsection{Compilazione e Avvio}
	\begin{lstlisting}[language=bash]
		# Avvio del Server
		$ ./server 8080
		
		# Avvio del Client
		$ ./client 127.0.0.1 8080
	\end{lstlisting}
	
	
	\subsection{Esempio di Interazione}
	Un tipico flusso di comunicazione tra Client (C) e Server (S):
	\begin{lstlisting}
		C: SEARCH Mario
		S: 200 OK 3331234567
		
		C: SEARCH Giuseppe
		S: 404 NOT_FOUND
		
		C: QUIT
		S: 200 BYE
	\end{lstlisting}
	
	\section{Conclusioni}
	Il progetto ha permesso di verificare l'efficacia della gestione delle socket in ambiente di rete e la semplicità di definizione di un protocollo applicativo su misura per servire richieste concorrenti.
	
\end{document}